% New checkpoint33 source: exact original E.61/E.65 frame comparison.
% Integral J is distinguished from C=J mod 3 throughout.
\statement{A.105}{Proved}{The original CM character is the product of the two actual frame signs}
Let \(\alpha=(-1,\phi)\in\widehat{\mathrm{GT}}_0\) be the actual
negative involution, with its canonical norm word \(u\), and put
\(\Sigma=\alpha c\), where \(c\) is the retained Galois complex
conjugation. Retain the ORIGINAL E.61/E.65 pure frames and coordinate:
\[
\begin{gathered}
 \Sigma(s)=s\phi=\beta_s s\beta_s^{-1},\qquad
 \Sigma(r)=\phi^{-1}r=A_r rA_r^{-1},\\
 A_r,\beta_s\in\widehat F,\qquad
 (\beta_s^{-1})^{\rm ab}=(\rho,\rho),\qquad \nu=\rho\bmod3.
\end{gathered}
\tag{A.105.1}
\]
Use the signed matrices \(X=I+2E_{12}\), \(Y=I-2E_{21}\) and
distinguish the integral matrix from its reduction:
\[
 J=\begin{pmatrix}0&-1\\1&-1\end{pmatrix},\qquad
 C=J\bmod3=XYX^{-1}\bmod3.
\tag{A.105.2}
\]
There are unique actual signs \(\epsilon_A,\epsilon_u\in\{1,-1\}\)
such that, in \(\mathrm{SL}_2(\mathbf F_3)\),
\[
 q_3(A_r)=\epsilon_A X^\nu C^{-\nu},\qquad
 q_3(u)=\epsilon_uY^{-\nu}X^\nu.
\tag{A.105.3}
\]
For the original character \(\Delta\) of E.65, including its
specified signed lift and rational three-torsion line,
\[
 \boxed{\,(-1)^{\Delta(\alpha)}
       =(-1)^{\delta(\Sigma)}=\epsilon_A\epsilon_u.\,}
\tag{A.105.4}
\]
If \(\alpha\in\widehat{\mathrm{GT}}\), put
\(d=1+2k\) for its actual even dihedral exponent \(k\). Then
\[
 \epsilon_u=d\bmod3,\qquad
 \boxed{\,\Delta(\alpha)=0
       \quad\Longleftrightarrow\quad \epsilon_A=d\bmod3.\,}
\tag{A.105.5}
\]
The equality on the right remains an open condition. In particular
this proposition does not assert vanishing of the original CM
character on the full Grothendieck--Teichmuller group.

\begin{proof}
\emph{A.105.1. The original finite inputs and their signs.}
The original E.61.3--4 and E.65.2--4 use
\[
\begin{gathered}
 U_0=\begin{pmatrix}1&1\\0&1\end{pmatrix},\qquad
 L_0=\begin{pmatrix}1&0\\-1&1\end{pmatrix},\\
 S_{65}=U_0L_0U_0,\qquad R_0=-U_0L_0=J,\\
 V=U_0^{8\rho},\qquad
 T=q_{12}(A_r)^{-1}V,\qquad
 W=q_{12}(\beta_s)^{-1}V.
\end{gathered}
\tag{A.105.6}
\]
These are signed matrices modulo twelve; \(q_{12}\) on the free
kernel sends \(x,y\) to the displayed \(X,Y\). The actual
multiplier-one finite covariance and parameter identity are
\[
\begin{gathered}
 q_{12}\Sigma=\operatorname{Ad}(V)q_{12},\\
 q_{12}(\phi)=
 \begin{pmatrix}1&-8\rho\\-8\rho&1+64\rho^2\end{pmatrix}.
\end{gathered}
\tag{A.105.7}
\]
These are the proved original E.61/E.65 identities for
\(\Sigma=(1,\phi)\), not identities for a substituted pair.
The source and complete referee comparison are bound in the
proof record below. In particular direct reduction gives
\[
 V_4=I,\qquad V_3=X^\nu,\qquad
 \Phi:=q_3(\phi)=Y^\nu X^{-\nu}.
\tag{A.105.8}
\]
The lift in (A.105.6) is the canonical even-power lift. Its cube
is \(I\) modulo twelve, whereas its negative has cube \(-I\).
No arbitrary sign is chosen independently for \(T\) or \(W\).

\emph{A.105.2. The actual signed norm phase.}
In this paragraph all matrices are modulo three. Put
\[
 S=\begin{pmatrix}0&-1\\1&0\end{pmatrix},\qquad
 D=\operatorname{diag}(1,-1),\qquad j=\operatorname{Ad}(D).
\]
Then \(j\) inverts \(X,Y\), while \(\operatorname{Ad}(S)\)
interchanges them. Since \(\alpha=\Sigma c\) and the actual
\(c\) inverts both free generators, (A.105.7) gives
\[
 q_3\alpha=\beta_\nu q_3,\qquad
 \beta_\nu=\operatorname{Inn}(X^\nu)j.
\tag{A.105.9}
\]
Thus this is descent of the actual automorphism, not a choice
of compatible finite matrices.

The derived subgroup of \(\mathrm{SL}_2(\mathbf F_3)\) is
\[
 Q_8=\langle S,R\rangle,\qquad
 R=Y^{-2}X^2=\begin{pmatrix}1&1\\1&2\end{pmatrix}.
\]
For completeness \(S^2=R^2=-I\), \(RSR^{-1}=S^{-1}\);
conjugation by \(X\) sends \(R\) to \(S\) and \(S\) to \(SR\).
These give a normal quaternion subgroup and the decomposition
\(Q_8\rtimes\langle X\rangle\) of the order-twenty-four matrix
group. Its commutators generate \(Q_8\), so its abelianization
is \(C_3\), with \(\chi(X)=\chi(Y)=1\).

The actual \(U=q_3(u)\) belongs to \(Q_8\), since \(u\in\widehat F'\).
Its canonical norm and duality give
\[
 U=\Phi\beta_\nu(U),\qquad SUS^{-1}=U^{-1}.
\tag{A.105.10}
\]
The fixed subgroup \(Q_8^{\beta_\nu}\) is \(\{I,-I\}\).
Indeed choose \(h\in\mathbf Z/3\) with \(2h=\nu\); conjugating
\(j\) by \(\operatorname{Inn}(X^h)\) gives \(\beta_\nu\).
The matrices in \(Q_8\) commuting with \(D\) are exactly the
diagonal matrices \(I,-I\).

The explicit element
\[
 U_\nu=Y^{-\nu}X^\nu
 =\begin{pmatrix}1&-\nu\\-\nu&1+\nu^2\end{pmatrix}
\tag{A.105.11}
\]
satisfies both equations (A.105.10): direct multiplication gives
\(\beta_\nu(U_\nu)=X^\nu Y^\nu X^\nu\), and
\(Y^{2\nu}=Y^{-\nu}\). If \(U_1,U_2\) satisfy the norm equation,
then \(U_1^{-1}U_2\) is \(\beta_\nu\)-fixed. Thus the actual
\(U\) has the unique form \(\epsilon_uU_\nu\), proving the
second assertion of (A.105.3). This is a classification of the
actual finite image; it does not posit two canonical profinite
norm solutions.

\emph{A.105.3. The original s-frame square.}
Work first in the actual free profinite product
\(\widehat B=C_2(s)\amalg C_3(r)\), with
\(\widehat F=\ker(\widehat B\to S_3)\). Write
\(\iota=\operatorname{Ad}(s)|_{\widehat F}\).
Canonical duality gives \((su)^2=1\), and its \(S_3\) image
is the nontrivial transposition. The finite-factor torsion and
centralizer theorems \cite[Theorems A' and B']{A_HR1985}
give a unique \(\gamma\in\widehat F\) with
\[
 su=\gamma s\gamma^{-1},\qquad
 u=\iota(\gamma)\gamma^{-1}.
\tag{A.105.12}
\]
Indeed an initial conjugator has image in
\(C_{S_3}(\bar s)=\langle\bar s\rangle\), and right multiplication
by \(s\) normalizes it into \(\widehat F\). Uniqueness follows
from \(C_{\widehat B}(s)=\langle s\rangle\) and
\(\widehat F\cap\langle s\rangle=1\).

The actual norm says \(\alpha(su)=su\). Applying \(\alpha\)
to (A.105.12), using the original frame \(\alpha(s)=\beta_s
s\beta_s^{-1}\), and then normalized-conjugator uniqueness gives
\[
 \alpha(\gamma)\beta_s=\gamma,\qquad
 \beta_s=\alpha(\gamma)^{-1}\gamma.
\tag{A.105.13}
\]
Since \(\alpha\) acts by \(-I\) on \(\widehat F^{\rm ab}\),
\((\beta_s)_{\rm ab}=2\gamma_{\rm ab}\). The condition
\(u\in\widehat F'\) makes \(\gamma_{\rm ab}\) diagonal.
Thus \(\beta_s\) has even diagonal abelianization and
\(q_4(\beta_s)=I\), because the original \(X_4,Y_4\) commute
and have order two.

Put \(B=q_3(\beta_s)\), \(G=q_3(\gamma)\). Applying the
SIGNED free-kernel map to the identities already inside
\(\widehat F\), rather than mapping \(s\) itself to \(S\), gives
\[
 B=\beta_\nu(G)^{-1}G,\quad
 U=SGS^{-1}G^{-1},\quad
 \Phi=SBS^{-1}B^{-1}.
\tag{A.105.14}
\]
Here \(S^2=-I\), so the last equation yields
\[
 BSB^{-1}=S\Phi=X^\nu SX^{-\nu}.
\]
The centralizer of \(S\) in the matrix group is \(\langle S\rangle\):
solving the commuting equation gives \(aI+bS\), with
\(a^2+b^2=1\), hence precisely \(I,-I,S,-S\).
Therefore \(B=X^\nu S^j\). Since
\(\chi(B)=2\chi(G)=\nu\), write \(G=X^{-\nu}H\) with \(H\in Q_8\).
Substitution in (A.105.14) gives
\[
 X^{-\nu}B=j(H)^{-1}H,\qquad
 U=Y^{-\nu}[S,H]X^\nu.
\tag{A.105.15}
\]
Every \(H\in Q_8\) is \(S^aR^b\), \(b\in\{0,1\}\).
The literal relations \(j(S)=S^{-1}\), \(j(R)=RS\) give
\[
 j(H)^{-1}H=
 \begin{cases}S^{2a},&b=0,\\S^{-2a-1},&b=1,\end{cases}
 \qquad
 \bigl(j(H)^{-1}H\bigr)^2=[S,H]=(-I)^b.
\tag{A.105.16}
\]
Comparison with (A.105.11) now proves
\[
 \bigl(B^{-1}X^\nu\bigr)^2=\epsilon_u I.
\tag{A.105.17}
\]
Inverting an element of \(\langle S\rangle\) does not change
its square. This retains the exact inverse-frame order in \(W\).

\emph{A.105.4. The original r-frame and its line.}
The original words \(x=(sr)^{-2}\), \(y=sxs\) show that
\(\rho_{\rm cyc}=\operatorname{Ad}(r)\) cycles \(x,y,z\).
From (A.105.1),
\[
 \phi=\rho_{\rm cyc}(A_r)A_r^{-1}.
\]
The matrix of \(\rho_{\rm cyc}-I\) on the free abelianization
has determinant three. Since \(\phi\in\widehat F'\), its
adjugate and the injectivity of multiplication by three on
\(\widehat{\mathbb Z}\) give \(A_r\in\widehat F'\).
This argument is made before reduction modulo three.

Modulo three, \(C=J\bmod3\) conjugates \(X\) to \(Y\), \(Y\)
to \(-Z\), and \(-Z\) to \(X\). The discrepancy from the actual
cyclic substitution is a central character, trivial on \(A_r\).
For \(A_3=q_3(A_r)\), therefore,
\[
 \Phi=CA_3C^{-1}A_3^{-1},\qquad
 A_3CA_3^{-1}=\Phi^{-1}C=X^\nu CX^{-\nu}.
\tag{A.105.18}
\]
The centralizer of \(C\) is
\(\langle C,-I\rangle\). One can solve this directly by
writing \(C=I+N\), \(N^2=0\), \(N\ne0\); its invertible
commuting determinant-one matrices are the six
\(\pm(I+tN)\). Thus \(A_3=X^\nu h\) with
\(h\in\langle C,-I\rangle\). Since \(A_3\in Q_8\) and
\(\chi(C)=1\), necessarily
\[
 A_3=\epsilon_A X^\nu C^{-\nu},\qquad
 T_3=A_3^{-1}X^\nu=\epsilon_A C^\nu.
\tag{A.105.19}
\]
This proves the first assertion of (A.105.3).
Also \(q_4(A_r)=I\), since the mod-four image is abelian.

\emph{A.105.5. Bind the named original characters.}
The original E.65 definition is
\[
\begin{gathered}
 \ell=\mathbf F_3(1,-1)=\ker(C-I),\\
 T_3|_\ell=(-1)^{\epsilon_3(T)}\,{\rm id}_\ell,\qquad
 \epsilon_2(S_{65}^j)=j\bmod2,\\
 \delta(\Sigma)=\epsilon_3(T)+\epsilon_2(W).
\end{gathered}
\tag{A.105.20}
\]
The outage matrix \(S\) used above is \(S_{65}^{-1}\); they
generate the same order-four group and define the same exponent
parity. Since \(C^\nu\) fixes \(\ell\) pointwise, (A.105.19)
gives \((-1)^{\epsilon_3(T)}=\epsilon_A\). Likewise
\(W_3^2=(-1)^{\epsilon_2(W)}I\), so (A.105.17) gives
\((-1)^{\epsilon_2(W)}=\epsilon_u\). Both corrected frames
are \(I\) modulo four, exactly as required by the original
global-sign normalization.

E.65's extension \(\Delta\) vanishes on the actual Galois
image. Hence \(\Delta(c)=0\), and
\(\Delta(\alpha)=\delta(\alpha c)=\delta(\Sigma)\).
This proves (A.105.4), with the original named characters.
For full-GT \(\alpha\), A.103 gives \(dU_{11}=1\).
Equation (A.105.11) gives \(U_{11}=\epsilon_u\) modulo three;
thus \(d\bmod3=\epsilon_u\), proving (A.105.5).
\end{proof}

The actual negative \(r\)-conjugator is
\[
 J_r=s\phi A_r=s\,\rho_{\rm cyc}(A_r),\qquad
 \alpha(J_r)J_r=1.
\]
It belongs to the non-pure \(s\)-coset. This equation does not
imply \(\alpha(A_r)=A_r^{-1}\); replacing \(J_r\) by \(A_r\)
would be a different, unjustified norm problem. The remaining
\(\epsilon_A=d\bmod3\) condition is not removed by the
all-congruence vanishing of \(H_d\).

{\small\raggedright
The full original-source binding, including the 327-line E.65
source, 267-line E.61 source and complete 252-line referee report,
with exact SHA-256 values and the two complete frame calculations,
is retained in
\path{work_DE/original_CM_source_binding33.md}.
The independent all-congruence audit is
\path{work_DE/all_signed_congruence_audit33.md}.
The integral matrix \(J\) is never identified with the integral
word \(XYX^{-1}\); only their reductions modulo three agree.
No old certificate was rerun, and this source does not assign a
value to the remaining CM parity.\par}
